\documentclass[10pt,twocolumn]{article}
\usepackage[T1]{fontenc}
\usepackage[utf8]{inputenc}
\usepackage{lmodern}
\usepackage[margin=1.8cm,top=2.4cm,bottom=2cm,columnsep=0.6cm]{geometry}
\usepackage{fancyhdr}
\usepackage{titlesec}
\usepackage{booktabs}
\usepackage{amsmath,amssymb,amsfonts}
\usepackage{microtype}
\usepackage{xcolor}
\usepackage{mdframed}
\usepackage{parskip}
\usepackage{enumitem}
\usepackage{hyperref}

\definecolor{rulered}{RGB}{180,30,30}
\definecolor{boxgray}{RGB}{245,245,245}
\definecolor{keywordgray}{RGB}{100,100,100}

\pagestyle{fancy}
\fancyhf{}
\fancyhead[L]{\textsc{\small Claude Mag}}
\fancyhead[C]{\small\textit{Machine Learning Research Digest}}
\fancyhead[R]{\small 2026-07-06}
\fancyfoot[C]{\small\thepage}
\renewcommand{\headrulewidth}{0.8pt}

\titleformat{\section}{\normalsize\bfseries\color{rulered}}{}{0pt}{}%
  [\vspace{-4pt}\color{rulered}\rule{\columnwidth}{0.4pt}]
\titlespacing{\section}{0pt}{8pt}{4pt}

\hypersetup{colorlinks=true,urlcolor=rulered,linkcolor=black}

\begin{document}
\setlength{\parindent}{0pt}
\setlength{\parskip}{4pt}

{\small\color{keywordgray}\textbf{Keywords:}
  reinforcement learning \textbullet{}
  maximum likelihood \textbullet{}
  LLM training \textbullet{}
  policy gradient}\par\vspace{4pt}

{\LARGE\bfseries\raggedright
  Maximum Likelihood Reinforcement Learning}\par\vspace{4pt}

{\large\itshape\raggedright
  REINFORCE Optimizes a Lower-Bound Approximation of True Likelihood ---
  MaxRL Fixes This with More Sampling Compute}\par\vspace{6pt}

{\small\textbf{By} Fahim Tajwar, Guanning Zeng, Yueer Zhou, Yuda Song,
  Daman Arora, Yiding Jiang, Jeff Schneider, Ruslan Salakhutdinov,
  Haiwen Feng \& Andrea Zanette (Carnegie Mellon University)}\par\vspace{2pt}

{\small\color{keywordgray}arXiv:2602.02710 \textbullet{}
  ICML 2026 Oral Presentation}
\par\vspace{2pt}\noindent{\color{rulered}\rule{\columnwidth}{0.6pt}}\vspace{6pt}

Standard reinforcement learning for LLMs does not maximize the likelihood
of correct rollouts --- it optimizes a Jensen lower bound that allows
probability to diffuse across mediocre samples. \textit{Maximum Likelihood
Reinforcement Learning} (MaxRL) introduces a compute-indexed family of
objectives that converge to exact maximum likelihood as more rollout samples
are drawn per query, achieving consistent improvements on math reasoning,
navigation, and large-scale Qwen3 training at no increase in gradient steps.

\begin{mdframed}[backgroundcolor=boxgray,linecolor=rulered,linewidth=1pt,
    innerleftmargin=6pt,innerrightmargin=6pt,
    innertopmargin=6pt,innerbottommargin=6pt]
\textbf{\small AT A GLANCE}\par\vspace{4pt}
\begin{description}[leftmargin=0pt,labelsep=4pt,itemsep=2pt,parsep=0pt,
    font=\small\bfseries]
  \item[Problem:] REINFORCE optimizes only a lower-order approximation of
    true maximum likelihood over correct rollouts in sampling-based RL settings.
  \item[Why it matters:] The objective gap causes probability mass to diffuse
    across mediocre rollouts rather than concentrating on correct answers.
  \item[Method:] Importance-weighted MLE estimator with $k$ samples per query,
    interpolating between REINFORCE ($k=1$) and exact MLE ($k\to\infty$).
  \item[Result:] Consistent improvement on maze navigation, GSM8K, and
    large-scale Qwen3 training; gains scale with $k$.
  \item[Evidence:] Toy classification verifies closeness to MLE; ablations
    confirm the importance-weighting is the key ingredient.
\end{description}
\end{mdframed}

\section{The Big Picture}

Training language models with reinforcement learning on binary correctness
signals --- did the generated code pass the tests? does the math proof check
out? --- is the dominant paradigm for post-training. The underlying objective
is maximum likelihood: given that a correct response exists, maximize the
probability the model assigns to it.

But the paper reveals a gap. REINFORCE, the workhorse of this paradigm,
does not maximize this likelihood. It optimizes a lower bound, one that is
tight only when the model already concentrates most probability on correct
responses. In practice, especially early in training, REINFORCE wastes gradient
signal on rolling updates that nudge mediocre rollouts upward rather than
concentrating on correct ones.

MaxRL closes this gap by drawing $k > 1$ samples per query and constructing
an importance-weighted objective whose gradient converges to the true MLE
gradient as $k$ grows. More sampling compute buys a closer approximation
to the ideal training signal.

\section{Why Existing Approaches Fall Short}

\textbf{REINFORCE:} With one sample per query and binary reward,
the REINFORCE gradient is:
\[
\nabla_\theta J_\text{REINFORCE} =
  \mathbb{E}_{y \sim \pi_\theta}\!\left[r(y) \nabla_\theta \log \pi_\theta(y|x)\right].
\]
This equals the gradient of $\mathbb{E}_x[\mathbb{E}_{y \sim \pi_\theta}[r(y) \log
\pi_\theta(y|x)]]$, a lower bound on $\mathbb{E}_x[\log \sum_{y \in Y^+}
\pi_\theta(y|x)]$ by Jensen's inequality. The gap is non-negligible when
$\pi_\theta$ spreads probability across many incorrect responses.

\textbf{Group Relative Policy Optimization (GRPO):} Draws $k$ samples but
uses them for reward normalization rather than MLE approximation; the
resulting objective is still a REINFORCE variant, not a MLE estimator.

\section{Core Mechanism}

MaxRL forms an importance-weighted estimate of the true MLE gradient using
$k$ rollouts per query. The weights rescale each rollout's contribution
proportionally to the model's current assignment of probability to that rollout
among all correct rollouts sampled:

\[
w_i = \frac{r(y_i)\,\pi_\theta(y_i|x)}
           {\displaystyle\sum_{j=1}^k r(y_j)\,\pi_\theta(y_j|x)}.
\]

The MaxRL objective is:
\[
J_\text{MaxRL}^{(k)}(\theta) =
  \mathbb{E}_{x,\,y_{1:k} \sim \pi_\theta}
  \!\left[\log \frac{1}{k} \sum_{i=1}^k r(y_i)\,\pi_\theta(y_i|x)\right],
\]
whose gradient is:
\[
\nabla_\theta J_\text{MaxRL}^{(k)} =
  \mathbb{E}\!\left[\sum_{i=1}^k w_i \,\nabla_\theta \log \pi_\theta(y_i|x)\right].
\]

At $k=1$ this reduces to REINFORCE; as $k \to \infty$ this converges to
the exact MLE gradient almost surely.

\section{Technical Formulation}

Let $Y^+$ be the set of correct responses for query $x$. The true
maximum likelihood objective is:
\[
J_\text{MLE}(\theta) = \mathbb{E}_{x}\!\left[\log \sum_{y \in Y^+} \pi_\theta(y|x)\right].
\]

\textbf{REINFORCE lower bound:} By Jensen's inequality applied to the
concave $\log$,
\[
J_\text{REINFORCE}(\theta) \leq J_\text{MLE}(\theta),
\]
with equality only when all probability mass is on a single correct response.

\textbf{MaxRL convergence:}
\[
J_\text{MaxRL}^{(k)}(\theta) \xrightarrow{\;k \to \infty\;} J_\text{MLE}(\theta)
\quad \text{a.s.}
\]
by the law of large numbers applied to the importance-weighted estimator.

The gap between MaxRL and MLE at finite $k$ is bounded by:
\[
|J_\text{MLE} - J_\text{MaxRL}^{(k)}| = O\!\left(k^{-1/2}\right)
\]
under mild regularity conditions on $\pi_\theta$.

\section{Learning or Inference Procedure}

\begin{enumerate}[itemsep=2pt,parsep=0pt]
  \item For each query $x$, sample $k$ rollouts $y_1, \ldots, y_k \sim \pi_\theta(\cdot|x)$.
  \item Evaluate reward $r(y_i) \in \{0,1\}$ for each rollout.
  \item Compute importance weights $w_i$ using current policy probabilities.
  \item Compute $\nabla_\theta J_\text{MaxRL}^{(k)}$ using the weighted
    gradient estimate.
  \item Update $\theta$ via Adam.
  \item At inference: standard autoregressive sampling (MaxRL changes
    only the training objective, not the inference procedure).
\end{enumerate}

\section{What the Guarantee Says}

MaxRL is a consistent estimator of the MLE gradient: as $k$ grows,
the training objective converges to exact maximum likelihood. The rate of
convergence is $O(k^{-1/2})$ in the objective value. No distributional
shift correction or importance weight clipping is required as long as
$\pi_\theta$ remains supported on $Y^+$.

\section{Experimental Findings}

\begin{itemize}[itemsep=2pt,parsep=0pt]
  \item \textbf{Toy classification:} MaxRL closely tracks exact MLE training
    signal; REINFORCE lags by a measurable margin throughout.
  \item \textbf{Maze navigation:} MaxRL achieves higher success rate than
    REINFORCE and GRPO at equivalent gradient steps.
  \item \textbf{GSM8K:} 3--5\% accuracy improvement over REINFORCE with
    the same compute; gains scale with $k$ up to $k = 16$.
  \item \textbf{Qwen3 (large-scale):} Consistent accuracy improvement;
    MaxRL with $k=8$ matches REINFORCE with $k=1$ at roughly half the
    gradient steps.
\end{itemize}

\section{Ablations and Interpretation}

\textbf{Effect of $k$:} Performance improves monotonically with $k$ up to
$k \approx 16$, beyond which gains plateau (model capacity becomes the
bottleneck, not objective quality).

\textbf{Importance weights vs.\ uniform:} Replacing $w_i$ with uniform
weights (equivalent to supervised fine-tuning on correct samples) degrades
performance, confirming that the reweighting term is the key ingredient,
not merely drawing more samples.

\textbf{Relationship to GRPO:} GRPO's reward normalization does not
approximate MLE; it is a variance-reduction technique within the REINFORCE
objective. MaxRL's improvement is orthogonal and additive to GRPO's tricks.

\section*{Reference}

Fahim Tajwar, Guanning Zeng, Yueer Zhou, Yuda Song, Daman Arora,
Yiding Jiang, Jeff Schneider, Ruslan Salakhutdinov, Haiwen Feng,
Andrea Zanette. \textbf{Maximum Likelihood Reinforcement Learning}.
arXiv:2602.02710, February 2026.
\url{https://arxiv.org/abs/2602.02710}

\end{document}
